\documentclass[11pt]{article}

\usepackage[margin=25mm]{geometry}

\usepackage{amsmath,amssymb,fontspec,unicode-math}

\IfFontExistsTF{DejaVu Serif}{\setmainfont{DejaVu Serif}}{\setmainfont{Latin Modern Roman}}

\IfFontExistsTF{Latin Modern Math}{\setmathfont{Latin Modern Math}}{\setmathfont{DejaVu Math TeX Gyre}}

\usepackage{xurl,hyperref}

\setlength{\parindent}{0pt}\setlength{\parskip}{6pt}

\begin{document}

\title{Reviving DSP for Advanced Theorem Proving in the Era of Reasoning Models}

\author{ELVIA preliminary AI analysis}\date{21 October 2026}\maketitle

Preliminary assessment: Minor.\par

Independent scientific review remains pending. Author adoption has not been requested.

Current success bookkeeping and hypothesis-fallback paths retain local defects. The authors’ historical totalized-logarithm domain issue is already corrected in the current benchmark statement.

\section*{Review scope}

Rapid current-version reassessment of every substantive concern in the supplied LaTeX report: the implicated paper passages and relevant author-released source paths were inspected, existing proposed corrections were briefly checked, and impact was calibrated to the demonstrated scope. No author-provided or historical verification program was executed.

\section*{Problem 1: The released success path does not enforce preservation of the requested theorem}

Source: Section 3.3, PDF pp. 4–5; Section 5.1, PDF p. 5; released Prove.run and verifier.

Assessment: supported; impact: minor

\subsection*{Source claim (paraphrase)}

Sections 3.3 and \(5.1\), PDF pp. 4–5, describe completing the requested proof and verifying it with Lean 4. Released dsp/prove.py, check\_proof, warns when the requested theorem is changed or commented out, but run can still mark an otherwise error-free, sorry-free returned program as finished; run\_direct returns the underlying completion result without enforcing target identity.

\subsection*{Analysis}

An unrelated declaration proving True can be error-free and sorry-free without proving a requested False statement. A target-mismatch warning that does not block the success branch therefore fails to enforce theorem identity. This is a controller-contract argument; no false Lean theorem or actual kernel acceptance of a false proposition is claimed.

Placing the exact requested target text inside a Lean block comment also passes the whitespace-substring check. Thus elevating the existing warning alone would not repair target preservation. These are conditional controller witnesses using a successful compiler-response interface, not claims of a newly compiled proof.

\subsection*{Finding and scope}

The current controller can accept a different completed declaration, but no actual target mismatch was found in the historical released-solution inventory and no benchmark score is revised.

\subsection*{Proposed correction}

Keep the target declaration immutable and restrict generation/repair to its proof body. Make acceptance depend on compiler-aware declaration name/type/assumption and admitted-axiom checks; a comment-stripping specimen guard is not a complete Lean parser.

A local correction or qualification is proposed where specified; author adoption is not established.

\section*{Problem 2: The unrestricted-hypothesis fallback is inactive by default and fails when enabled}

Source: Appendix A.1, PDF p. 21; released fallback send\_request.

Assessment: supported; impact: minor

\subsection*{Source claim (paraphrase)}

Appendix A.1, PDF p. 21, describes both hinted-hypothesis and unrestricted-context proof search. In the released proving interface, run and run\_direct leave the strict retry disabled by default; its enabled branch calls send\_request without a required argument.

\subsection*{Analysis}

Normal calls leave the strict flag false, so the unrestricted retry is not reached. Enabling it exposes a missing required send\_request argument and fails before the request. Both mechanisms are version-bounded control-flow/interface defects; their benchmark incidence is unmeasured.

\subsection*{Finding and scope}

The current fallback is unreachable under default calls and raises a missing-argument error when enabled; search coverage changes, without invalidating accepted Lean proofs.

\subsection*{Proposed correction}

Pass both proof blocks and the current bfs\_list to the retry, and explicitly enable the intended fallback policy. This restores the local interface, without establishing an improved theorem-proving success rate.

A local correction or qualification is proposed where specified; author adoption is not established.

\section*{Problem 3: Reproduction of the authors’ known logarithm-domain issue in miniF2F}

Source: Appendix F, PDF pp. 27–28 (known nested-logarithm contest formalization).

Assessment: not\_issue; impact: no\_issue\_found

\subsection*{Source claim (paraphrase)}

Appendix F, PDF pp. 27–28, already diagnoses and corrects the original miniF2F nested-logarithm formalization. The inspected author dataset contains the strengthened domain condition. No discovery credit is claimed.

\subsection*{Analysis}

Under Lean's totalized Real.logb convention, \(n=1\) makes both nested logarithms zero and satisfies the original formal equation, while its decimal digit sum is \(1\), not \(13\). The informal contest problem's positive logarithm domains exclude this endpoint. This reproduces the authors' disclosed formalization issue.

The current author dataset explicitly requires \(1<n\), excluding this witness. On that domain, writing \(t=\log_2 n>0\) reduces the equation to \(\log_2t-2=(\log_2t-1)/2\), hence \(n=256\) and digit sum \(13\).

\subsection*{Finding and scope}

The author-disclosed historical domain error is corrected in the current dataset. It is retained as history and supplies no current issue severity.

\section*{Limits}

No author-provided program, training procedure, model inference or wet-lab experiment was executed in this review.\par

Historical execution results are retained only with that attribution; a source-level defect alone does not identify the program or configuration used for the published experiments.\par

No author contact, adoption, independent human scientific review or novelty determination is established.\par

The findings apply only to the identified claims, source versions and stated assumptions; they do not certify or invalidate the entire paper.\par

\section*{Source identity}

NeurIPS 2025 proceedings (39 pages); current arXiv:2506.11487v1 compared

\url{https://papers.neurips.cc/paper_files/paper/2025/file/6b44ee74539ea77d6a0d50d468724371-Paper-Conference.pdf}

Reviewed PDF SHA-256: \nolinkurl{bf9786826ebead7053526acf960a497244d09d9f80a6cef2388741304ae1233d}

The exact inspected current paper PDF is identified by its SHA-256. The current formal conference version includes expanded benchmark-correction discussion. The latest arXiv channel remains the earlier June manuscript with different pagination. Current DSP-Plus head eb983208b78b5ac4a7774d5500cc2086a2171027 retains controller/fallback defects but already has the strengthened logarithm domain. An exact PDF acquisition timestamp is not established in this reassessment; source-currentness checks have their own recorded timestamps and file modification times were not substituted.

\end{document}
